\scanpage{22}\mspage{21}
Or si $B'\in G$ est tel que $A'^\circ\subset B'^\circ$, on a \deleted{$A\subset B'$} (hyp. de b) $B'\subset \deleted{A}B$, donc $A'^\circ\subset B$, OK.

Prouvons que $B$ est le plus petit él. de $G^A$. Soit $B'\in G^A$ i.e. $B'\in G$, $A\subset B'$, prouvons $B\subset B'$ ou encore $B^\circ\deleted{\unsure{\subset}}\cap B'\ne\varnothing$ \deleted{Comme}. Comme
$A^\circ\subset B^\circ$ et $A^\circ\subset A\subset B'$, on a $\underset{\substack{\mathrel{\ooalign{\hfil$\scriptstyle\shortparallel$\hfil\cr\hfil$\scriptstyle-$\hfil\cr}}\\ \varnothing}}{A^\circ}\subset B^\circ\cap B'$, OK.

Prouvons enfin que $A\in F$, $B\in G$ implique que $A\cap B$ est $F$ fermé dans $|F|$. Par b) on voit que $\underset{\substack{\shortparallel\\ \bigcup_{A\in F}A^\circ}}{|F|}\subset|G|$, et que
la relation d'équiv. $\unsure{\mathcal R|R}_F$ sur $|F|$ est plus fine que celle induite par $\mathcal R_G$ sur $|G|$. Cela implique que $A\cap B$ est saturée pour la relation d'équiv. $\deleted{\unsure{F}}\mathcal R_F$ (puisque $B$ l'est pour $\mathcal R_G$). Donc
\[ A\cap B=\bigcup_{\substack{C\in F\\ \text{t.q. } C^\circ\subset A\cap B}} C^\circ . \]
Je dis que \inserted{dans} cette réunion, \deleted{est} on peut remplacer $C^\circ$ par $C$ i.e.
\[ C\in F,\quad C^\circ\subset A\cap B\quad\Longrightarrow\quad C\subset A\cap B. \]
Comme $C^\circ\subset A$, on a $C\deleted{\unsure{{}^\circ}}\subset A$. Reste : prouver $C\subset B$. On va donner un complément \unsure{immédiat|inédit} à l'énoncé de la prop.
% slanted
\margin{Ou encore : ($\mathrm{b}''$) \unsure{On|Ou|On a} : $|F|\subset|G|$, la relation de préordre sur $|F|$ définie par $F$ est plus fine que celle induite par la relation d'ordre de $|G|$.}

\uline{Corollaire} La condition b) équivaut à
$\mathrm{b}'$) $\forall\ A\in F,\ \exists\ B\in G$ tel que $A^\circ\subset B^\circ$, et d'autre part $A\in F$, $B\in G$, $A^\circ\subset B$ implique $A\subset B$.
